\documentclass[11pt]{article}
% Shared preamble for all solution files.
\usepackage[margin=1.1in]{geometry}
\usepackage{enumitem}
\usepackage{tikz}
\usepackage[hidelinks]{hyperref}
\usepackage{amsmath,amssymb,amsthm}
\newcommand{\fa}{\mathfrak a}
\newcommand{\fb}{\mathfrak b}
\newcommand{\fc}{\mathfrak c}
\newcommand{\fp}{\mathfrak p}
\newcommand{\fq}{\mathfrak q}
\newcommand{\fm}{\mathfrak m}
\newcommand{\fn}{\mathfrak n}
\newcommand{\fr}{\mathfrak r}
\newcommand{\Nil}{\mathfrak N}
\newcommand{\Jac}{\mathfrak R}
\newcommand{\rad}{r}
\newcommand{\Spec}{\operatorname{Spec}}
\newcommand{\Max}{\operatorname{Max}}
\newcommand{\Ann}{\operatorname{Ann}}
\newcommand{\Ker}{\operatorname{Ker}}
\newcommand{\Z}{\mathbf Z}
\newcommand{\Q}{\mathbf Q}
\newcommand{\R}{\mathbf R}
\newcommand{\C}{\mathbf C}
\newcommand{\ci}{\mathrm{c}}
\newcommand{\ex}{\mathrm{e}}

\setlength{\parindent}{0pt}
\setlength{\parskip}{0.55em}
% \exercise{label}{statement}  -- followed by the solution text
\newcommand{\exercise}[2]{\par\bigskip\noindent{\large\textbf{Exercise #1}}\par\smallskip
  \begin{quote}\small\textbf{Problem.}\ #2\end{quote}\par}
\newcommand{\pf}{\noindent\textbf{Solution.}\ }

\title{Atiyah--Macdonald, Chapter 11: Dimension Theory\\[2pt]\large Solutions to the exercises}
\author{}
\date{}
\begin{document}
\maketitle
\noindent\textbf{Conventions.} Rings are commutative with $1$; $\operatorname{ht}$ denotes height of a prime ideal and $\dim$ the Krull dimension. From the text I use: Krull's principal ideal theorem in the form 11.16 (in a Noetherian ring every minimal prime ideal of $(x_1,\dots,x_r)$ has height $\le r$); 11.12 (heights are finite in a Noetherian ring); 11.13 (a Noetherian local ring of dimension $d$ has an $\fm$-primary ideal generated by $d$ elements); 11.14 (dimension theorem), together with the example after it ($\dim k[x_1,\dots,x_n]_{(x_1,\dots,x_n)}=n$); 11.18 ($\dim A/(x)=\dim A-1$ for a non-zero-divisor $x\in\fm$); 11.20 and 11.21; 11.22 (regularity: $\dim_k\fm/\fm^2=\dim A$); 11.26 (integral extension of integrally closed domains preserves local dimension). Also the Normalization Lemma (Ch.~5, Ex.~16), 5.8--5.10 (maximal ideals, lying over), 10.24 (finite generation of complete modules), Ch.~3 Ex.~21 and Ch.~4 Ex.~7 (the ideals $\fp[x]$).

This chapter has no numbered in-text exercises; there are 7 end-of-chapter exercises. Exercise 4 (Nagata) asks us to verify several assertions made in its statement; each is proved below. Exercise 5 refers to the Grothendieck group $K(A_0)$ of Ch.~7, Ex.~26, which I recall in the statement.

%=====================================================================
\section*{End-of-chapter exercises}

%---------------------------------------------------------------- 1
\exercise{1}{Let $f\in k[x_1,\dots,x_n]$ be an irreducible polynomial over an algebraically closed field $k$. A point $P$ on the variety $f(x)=0$ is \emph{non-singular} $\iff$ not all the partial derivatives $\partial f/\partial x_i$ vanish at $P$. Let $A=k[x_1,\dots,x_n]/(f)$, and let $\fm$ be the maximal ideal of $A$ corresponding to the point $P$. Prove that $P$ is non-singular $\iff A_\fm$ is a regular local ring. [By (11.18) we have $\dim A_\fm=n-1$. Now $\fm/\fm^2\cong(x_1,\dots,x_n)/\bigl((x_1,\dots,x_n)^2+(f)\bigr)$ and has dimension $n-1$ if and only if $f\notin(x_1,\dots,x_n)^2$.]}
\pf
Write $P=(a_1,\dots,a_n)$, $B=k[x_1,\dots,x_n]$, $M=(x_1-a_1,\dots,x_n-a_n)$ (the maximal ideal of $P$ in $B$). Since $P$ lies on the variety, $f\in M$, and $f\ne0$. The ring $A=B/(f)$ is a domain ($f$ irreducible, $B$ a UFD), and $\fm=M/(f)$.

\emph{Dimension.} $B_M$ is a Noetherian local domain of dimension $n$ (translate $P$ to the origin and use the example after 11.14). $f\in M B_M$ is a non-zero element of a domain, hence a non-zero-divisor in the maximal ideal. By 11.18, $\dim A_\fm=\dim B_M/(f)=n-1$.

\emph{Cotangent space.} $\fm/\fm^2$ is a vector space over $k=A/\fm$ killed by $\fm$, so it is unchanged by localization at $\fm$: $\fm A_\fm/\fm^2A_\fm\cong\fm/\fm^2$. Moreover $\fm^2=(M^2+(f))/(f)$, so
\[
\fm/\fm^2\cong M/(M^2+(f))=\bigl(M/M^2\bigr)\big/\bigl(\text{image of }f\bigr).
\]
$M/M^2$ has $k$-basis the classes of $x_i-a_i$ ($1\le i\le n$), so $\dim_kM/M^2=n$. Expand $f$ in powers of $y_i=x_i-a_i$ (Taylor expansion at $P$, valid for formal derivatives in every characteristic): $f=f(P)+\sum_i\frac{\partial f}{\partial x_i}(P)\,y_i+(\text{terms in }M^2)$, with $f(P)=0$. Hence the image of $f$ in $M/M^2$ is $\sum_i\frac{\partial f}{\partial x_i}(P)\,\bar y_i$, which is non-zero iff some partial derivative does not vanish at $P$. So
\[
\dim_k\fm/\fm^2=\begin{cases}n-1&\text{if }P\text{ is non-singular},\\ n&\text{if }P\text{ is singular.}\end{cases}
\]
By 11.22, $A_\fm$ is regular iff $\dim_k\fm/\fm^2=\dim A_\fm=n-1$, i.e.\ iff $P$ is non-singular. \hfill$\square$

%---------------------------------------------------------------- 2
\exercise{2}{In (11.21) assume that $A$ is complete. Prove that the homomorphism $k[[t_1,\dots,t_d]]\to A$ given by $t_i\mapsto x_i$ ($1\le i\le d$) is injective and that $A$ is a finitely-generated module over $k[[t_1,\dots,t_d]]$. [Use (10.24).]}
\pf
Recall the setting of 11.21: $A$ is a Noetherian local ring with maximal ideal $\fm$, $k\subseteq A$ a field mapping isomorphically onto $A/\fm$, $x_1,\dots,x_d$ a system of parameters ($d=\dim A$), and $\fq=(x_1,\dots,x_d)$ is $\fm$-primary. Put $B=k[[t_1,\dots,t_d]]$.

\emph{The homomorphism.} Since $\fm^r\subseteq\fq\subseteq\fm$ for some $r$, the $\fq$-adic and $\fm$-adic topologies on $A$ agree, and $A$ is complete for them. For $g=\sum c_\alpha t^\alpha\in B$ ($c_\alpha\in k$) the series $\sum c_\alpha x^\alpha$ converges in $A$ ($x^\alpha\in\fq^{|\alpha|}$), and $\phi(g)=\sum c_\alpha x^\alpha$ defines a $k$-algebra homomorphism $\phi:B\to A$ with $\phi((t)^n)\subseteq\fq^n$, where $(t)=(t_1,\dots,t_d)$.

\emph{Step 1: $G_{(t)}(B)=k[t_1,\dots,t_d]\to G_\fq(A)$, $t_i\mapsto\bar x_i$, is injective and makes $G_\fq(A)$ a finitely generated module.} This is the map induced by $\phi$ on associated graded rings. \emph{Injective:} the map is graded, so it suffices to show that no non-zero homogeneous $h\in k[t]$ of degree $s$ is in its kernel, i.e.\ that $h(x)\in\fq^{s+1}$ forces $h=0$. By 11.20, if $h(x)\in\fq^{s+1}$ then all coefficients of $h$ lie in $\fm$; but they lie in $k$ and $k\cap\fm=0$ ($k\to A/\fm$ is injective). So $h=0$. \emph{Finite generation:} $A/\fq$ has finite length (it is Artin), so it is a finite-dimensional $k$-vector space (every composition factor is $A/\fm\cong k$), say spanned by $e_1,\dots,e_N$. $G_\fq(A)_n=\fq^n/\fq^{n+1}$ is spanned over $A/\fq$ by the classes of the monomials of degree $n$ in $x_i$, hence over $k$ by $e_j\cdot\bar x^\alpha$. So the $e_j\in G_\fq(A)_0$ generate $G_\fq(A)$ as a $k[t]$-module.

\emph{Finite generation of $A$ over $B$.} Regard $A$ as a $B$-module via $\phi$ with the filtration $A_n=\fq^n$; then $(t)A_n\subseteq A_{n+1}$, so this is a $(t)$-filtration. $B$ is complete in the $(t)$-adic topology, $A$ is Hausdorff in the filtration topology ($\bigcap\fq^n=0$ by Krull's theorem, 10.17, as $A$ is Noetherian local), and $G(A)$ is a finitely generated $G(B)$-module by Step 1. By 10.24, $A$ is a finitely generated $B$-module.

\emph{Injectivity of $\phi$.} Let $0\ne g\in B$ with $\phi(g)=0$. Write $g=g_s+g_{s+1}+\cdots$ with $g_j$ homogeneous of degree $j$ and $g_s\ne0$. Every monomial $x^\alpha$ with $|\alpha|>s$ can be written $x^\beta x^{\alpha-\beta}$ with $|\beta|=s+1$ (choose $\beta$ for each $\alpha$), so grouping the convergent series $\sum_{|\alpha|>s}c_\alpha x^\alpha=\sum_{|\beta|=s+1}x^\beta u_\beta$ with $u_\beta\in A$; this lies in $\fq^{s+1}$. Then $0=\phi(g)=g_s(x)+(\text{element of }\fq^{s+1})$ gives $g_s(x)\in\fq^{s+1}$, hence $g_s=0$ by Step 1 (again 11.20), a contradiction. So $\phi$ is injective. \hfill$\square$

%---------------------------------------------------------------- 3
\exercise{3}{Extend (11.25) to non-algebraically-closed fields. [If $\bar k$ is the algebraic closure of $k$, then $\bar k[x_1,\dots,x_n]$ is integral over $k[x_1,\dots,x_n]$.]}
\pf
\emph{Statement.} Let $k$ be \emph{any} field, $\fp$ a prime ideal of $k[x_1,\dots,x_n]$, $A=k[x_1,\dots,x_n]/\fp$, $K$ its field of fractions and $d=\operatorname{tr.deg}_kK$. Then $\dim A_\fm=d$ for every maximal ideal $\fm$ of $A$; consequently $\dim A=d$. (For $k$ algebraically closed, maximal ideals are the points of the variety, and this is 11.25 and 11.27.)

\emph{Step 1: for $B=k[y_1,\dots,y_d]$ (independent indeterminates) and every maximal ideal $\fn$ of $B$, $\dim B_\fn=d$.} Let $\bar k$ be an algebraic closure of $k$ and $\bar B=\bar k[y_1,\dots,y_d]\supseteq B$. Every element of $\bar k$ is algebraic over $k$, hence integral over $B$; together with the $y_i\in B$ this shows $\bar B$ is integral over $B$ (the integral closure of $B$ in $\bar B$ is a subring containing $\bar k$ and $B$). $B$ is a UFD, hence integrally closed (5.12). By lying over (5.10) there is a prime $\bar\fm$ of $\bar B$ with $\bar\fm\cap B=\fn$, and $\bar\fm$ is maximal by 5.8. By Lemma 11.26, $\dim B_\fn=\dim\bar B_{\bar\fm}$. As $\bar k$ is algebraically closed, $\bar\fm=(y_1-c_1,\dots,y_d-c_d)$ for some $c_i\in\bar k$ (Nullstellensatz, 7.9), and after translating $c$ to the origin the example after 11.14 gives $\dim\bar B_{\bar\fm}=d$.

\emph{Step 2: the general case.} By the Normalization Lemma (Ch.~5, Ex.~16; the substitution $x_i\mapsto x_i-x_n^{N^i}$ works over any field) there are $y_1,\dots,y_r\in A$ algebraically independent over $k$ with $A$ integral over $B=k[y_1,\dots,y_r]$. Then $K$ is algebraic over $k(y_1,\dots,y_r)$, so $r=\operatorname{tr.deg}_kK=d$. $B$ is integrally closed (UFD) and $A$ is an integral domain integral over $B$. For a maximal ideal $\fm$ of $A$, $\fn=\fm\cap B$ is maximal (5.8) and Lemma 11.26 gives $\dim A_\fm=\dim B_\fn=d$ by Step 1. Hence $\dim A=\sup_\fm\dim A_\fm=d$ too. \hfill$\square$

%---------------------------------------------------------------- 4
\exercise{4}{(An example of a Noetherian domain of infinite dimension, Nagata.) Let $k$ be a field and let $A=k[x_1,x_2,\dots,x_n,\dots]$ be a polynomial ring over $k$ in a countably infinite set of indeterminates. Let $m_1,m_2,\dots$ be an increasing sequence of positive integers such that $m_{i+1}-m_i>m_i-m_{i-1}$ for all $i>1$. Let $\fp_i=(x_{m_i+1},\dots,x_{m_{i+1}})$ and let $S$ be the complement in $A$ of the union of the ideals $\fp_i$. Each $\fp_i$ is a prime ideal and therefore the set $S$ is multiplicatively closed. The ring $S^{-1}A$ is Noetherian by Chapter 7, Exercise 9. Each $S^{-1}\fp_i$ has height equal to $m_{i+1}-m_i$, hence $\dim S^{-1}A=\infty$.}
\pf
Put $r_i=m_{i+1}-m_i$. Call $X_i=\{x_{m_i+1},\dots,x_{m_{i+1}}\}$ the $i$-th block of variables; the blocks are pairwise disjoint.

\emph{$\fp_i$ is prime and $S$ is multiplicatively closed.} $A/\fp_i$ is the polynomial ring in all variables outside $X_i$, a domain. The complement of a union of primes is multiplicatively closed ($1\in S$). $A$ is a domain, so $S^{-1}A$ is a domain.

\emph{(a) A non-zero $a\in A$ lies in only finitely many $\fp_i$.} A polynomial involves finitely many variables. If no variable of $X_i$ occurs in $a$, then $a\in\fp_i$ would force $a=0$ (a polynomial lies in $\fp_i=(X_i)$ iff each of its monomials contains a variable of $X_i$).

\emph{(b) Every prime $\fq$ of $A$ with $\fq\cap S=\emptyset$ (i.e.\ $\fq\subseteq\bigcup\fp_i$) is contained in some $\fp_i$.} Suppose not; we construct an element of $\fq$ outside every $\fp_j$, contradicting $\fq\subseteq\bigcup\fp_j$. Clearly $\fq\ne0$ (as $0\subseteq\fp_1$); take $0\ne x\in\fq$ and let $I=\{i:x\in\fp_i\}$, finite by (a); $I\ne\emptyset$ since $x\notin S$. As $\fq\not\subseteq\fp_i$ for $i\in I$, prime avoidance (1.11) gives $z\in\fq$ with $z\notin\fp_i$ for all $i\in I$. Choose a variable $x_N$ not occurring in $x$ and not belonging to any block $X_i$ with $i\in I$ (possible: infinitely many variables). Put $u=x+zx_N\in\fq$.
\begin{itemize}[leftmargin=*]
\item For $i\in I$: $x\in\fp_i$ while $z,x_N\notin\fp_i$, so $zx_N\notin\fp_i$ (prime) and $u\notin\fp_i$.
\item For $j\notin I$: $x\notin\fp_j$, so $x$ has a monomial $\mu$ containing no variable of $X_j$. No monomial of $zx_N$ equals $\mu$, because every monomial of $zx_N$ is divisible by $x_N$ while $\mu$ is not. So $\mu$ occurs in $u$ with its coefficient unchanged, and $\mu$ has no variable of $X_j$; hence $u\notin\fp_j$.
\end{itemize}
Thus $u\in\fq\cap S$, a contradiction. So $\fq\subseteq\fp_i$ for some $i$.

\emph{(c) Maximal ideals of $R=S^{-1}A$.} Primes of $R$ are $S^{-1}\fq$ with $\fq$ a prime of $A$ disjoint from $S$ (3.11), inclusion-preserving. By (b) each such $\fq$ lies in some $\fp_i$; and $\fp_i\subseteq\fp_j$ forces $i=j$ (the blocks are disjoint, so $x\in X_i$ is not in $\fp_j$ for $j\ne i$). Hence the maximal ideals of $R$ are exactly the $\fn_i=S^{-1}\fp_i$.

\emph{(d) $R_{\fn_i}$ is Noetherian and $\dim R_{\fn_i}=r_i$.} Since $S\subseteq A\setminus\fp_i$, $R_{\fn_i}\cong A_{\fp_i}$ (Ch.~3, Ex.~3). Let $T=k[x_j:x_j\notin X_i]\setminus\{0\}\subseteq A\setminus\fp_i$ and let $F$ be the field of fractions of $k[x_j:x_j\notin X_i]$. Then $T^{-1}A=F[X_i]$ is the polynomial ring in $r_i$ variables over $F$, and $A_{\fp_i}=(T^{-1}A)_{T^{-1}\fp_i}=F[X_i]_{(X_i)}$, the localization of a polynomial ring at the ideal generated by all its variables. This is a Noetherian local ring of dimension $r_i$ (example after 11.14).

\emph{(e) $R$ is Noetherian.} By Ch.~7, Ex.~9 (a ring in which every localization at a maximal ideal is Noetherian and every non-zero element lies in only finitely many maximal ideals is Noetherian) it suffices to check these two conditions. The first is (d). For the second let $0\ne a/s\in R$; then $a\ne0$ and $a/s\in\fn_i\iff a\in\fp_i$ (as $\fp_i\cap S=\emptyset$), which holds for only finitely many $i$ by (a).

\emph{(f) Dimension.} $\operatorname{ht}\fn_i=\dim R_{\fn_i}=r_i=m_{i+1}-m_i$ by (d). The hypothesis $m_{i+1}-m_i>m_i-m_{i-1}$ says the $r_i$ form a strictly increasing sequence of integers, so $r_i\to\infty$. Therefore $\dim S^{-1}A\ge\sup_ir_i=\infty$: $S^{-1}A$ is a Noetherian domain of infinite Krull dimension. (Each local ring $R_{\fn_i}$ has finite dimension, in agreement with Cor.~11.11, but these dimensions are unbounded.) \hfill$\square$

%---------------------------------------------------------------- 5
\exercise{5}{Reformulate (11.1) in terms of the Grothendieck group $K(A_0)$ (Chapter 7, Exercise 26).

(Recall: for a Noetherian ring $A_0$, $K(A_0)$ is the quotient of the free abelian group on isomorphism classes $(M)$ of finitely generated $A_0$-modules by the subgroup generated by $(M')-(M)+(M'')$ for all short exact sequences $0\to M'\to M\to M''\to0$; $\gamma(M)$ denotes the class of $M$. By Ch.~7, Ex.~26(i), $\gamma$ is the \emph{universal} additive function: for every additive function $\lambda$ on finitely generated $A_0$-modules with values in an abelian group $G$ there is a unique homomorphism $\lambda_0:K(A_0)\to G$ with $\lambda=\lambda_0\circ\gamma$.)}
\pf
\emph{Reformulation.} Let $A=\bigoplus_{n\ge0}A_n$ be a Noetherian graded ring (so $A_0$ is Noetherian by 10.7), generated over $A_0$ by homogeneous elements $x_1,\dots,x_s$ of degrees $k_1,\dots,k_s>0$, and let $M=\bigoplus M_n$ be a finitely generated graded $A$-module. Each $M_n$ is a finitely generated $A_0$-module, so its class $\gamma(M_n)\in K(A_0)$ is defined. Define the \emph{universal Poincaré series}
\[
P_\gamma(M,t)=\sum_{n=0}^\infty\gamma(M_n)\,t^n\ \in K(A_0)[[t]].
\]
\textbf{Theorem 11.1$'$.} $P_\gamma(M,t)=\dfrac{f(t)}{\prod_{i=1}^s(1-t^{k_i})}$ with $f(t)\in K(A_0)[t]$; that is, $\prod_{i=1}^s(1-t^{k_i})\cdot P_\gamma(M,t)=f(t)$ in $K(A_0)[[t]]$. \qquad (For every additive function $\lambda$ with values in $\Z$ (or any abelian group), $P(M,t)=\sum\lambda(M_n)t^n$ is the image of $P_\gamma(M,t)$ under the coefficientwise map $\lambda_0:K(A_0)[[t]]\to\Z[[t]]$, so Theorem 11.1 follows by applying $\lambda_0$ to $f$.)

\emph{Proof} (the proof of 11.1 verbatim, with $\gamma$ in place of $\lambda$; it only uses that $\lambda$ is additive and takes values in an abelian group). Induction on $s$. If $s=0$ then $A=A_0$; $M$ is generated by finitely many homogeneous elements of degrees $\le r$, so $M_n=0$ for $n>r$ and $P_\gamma(M,t)$ is a polynomial.

Let $s>0$ and assume the result for rings generated by $s-1$ elements. Multiplication by $x_s$ is an $A_0$-linear map $M_n\to M_{n+k_s}$; let $K_n=\Ker$ and $L_{n+k_s}=\Coker$, giving exact sequences of finitely generated $A_0$-modules
\[
0\to K_n\to M_n\xrightarrow{x_s}M_{n+k_s}\to L_{n+k_s}\to0 .
\]
Put $K=\bigoplus K_n$, $L=\bigoplus L_n$ (with $L_n=M_n$ for $n<k_s$). $K$ is a submodule and $L=M/x_sM$ a quotient of $M$, so they are finitely generated $A$-modules (as $A$ is Noetherian), and both are annihilated by $x_s$; they are therefore finitely generated graded modules over the Noetherian graded ring $A/(x_s)$, which is generated over $A_0$ by the images of $x_1,\dots,x_{s-1}$. Since $\gamma$ is additive, (2.11) gives $\gamma(K_n)-\gamma(M_n)+\gamma(M_{n+k_s})-\gamma(L_{n+k_s})=0$. Multiply by $t^{n+k_s}$ and sum over $n\ge0$:
\[
t^{k_s}P_\gamma(K,t)-t^{k_s}P_\gamma(M,t)+\bigl(P_\gamma(M,t)-g(t)\bigr)-\bigl(P_\gamma(L,t)-h(t)\bigr)=0,
\]
where $g=\sum_{j<k_s}\gamma(M_j)t^j$ and $h=\sum_{j<k_s}\gamma(L_j)t^j$ are polynomials. Hence $(1-t^{k_s})P_\gamma(M,t)=P_\gamma(L,t)-t^{k_s}P_\gamma(K,t)+g(t)-h(t)$. By the inductive hypothesis $P_\gamma(L,t)$ and $P_\gamma(K,t)$ have the form (polynomial)$/\prod_{i<s}(1-t^{k_i})$, and the result follows. \hfill$\square$

%---------------------------------------------------------------- 6
\exercise{6}{Let $A$ be a ring (not necessarily Noetherian). Prove that
\[
1+\dim A\le\dim A[x]\le1+2\dim A.
\]
[Let $f:A\to A[x]$ be the embedding and consider the fiber of $f^*:\Spec(A[x])\to\Spec(A)$ over a prime ideal $\fp$ of $A$. This fiber can be identified with the spectrum of $k\otimes_AA[x]\cong k[x]$, where $k$ is the residue field at $\fp$ (Chapter 3, Exercise 21), and $\dim k[x]=1$. Now use Exercise 7(ii) of Chapter 4.]}
\pf
If $\dim A=\infty$ both inequalities say $\dim A[x]=\infty$, which follows from the lower bound below; so assume the dimension finite where needed.

\emph{Lower bound.} Let $\fp_0\subsetneq\cdots\subsetneq\fp_r$ be a chain of primes of $A$. By Ch.~4, Ex.~7(ii) the ideals $\fp_i[x]$ are prime in $A[x]$, and they form a strictly increasing chain (compare constant polynomials). Moreover $\fp_r[x]\subsetneq\fp_r[x]+(x)$, which is prime since $A[x]/(\fp_r[x]+(x))\cong A/\fp_r$ is a domain. This is a chain of length $r+1$, so $\dim A[x]\ge1+\dim A$.

\emph{Upper bound.} Let $Q_0\subsetneq\cdots\subsetneq Q_m$ be a chain of primes of $A[x]$ and $\fp_i=Q_i\cap A$, so $\fp_0\subseteq\cdots\subseteq\fp_m$. Fix a prime $\fp$ of $A$ and consider the primes $Q$ of $A[x]$ with $Q\cap A=\fp$. They contain $\fp[x]$, so correspond to primes of $(A/\fp)[x]$ whose contraction to $A/\fp$ is $0$, i.e.\ which do not meet $S=(A/\fp)\setminus\{0\}$; by 3.11(iv) these correspond, inclusion-preservingly, to primes of $S^{-1}((A/\fp)[x])=k(\fp)[x]$ (this is the identification of the fibre of Ch.~3, Ex.~21(iv)). Now $k(\fp)[x]$ is a principal ideal domain of dimension $1$, so it has no chain of three distinct primes. Hence at most two of the $Q_i$ have the same contraction $\fp$. The distinct values among $\fp_0\subseteq\cdots\subseteq\fp_m$ form a strictly increasing chain of primes of $A$, so there are at most $\dim A+1$ of them. Therefore $m+1\le2(\dim A+1)$, i.e.\ $m\le1+2\dim A$. Hence $\dim A[x]\le1+2\dim A$. \hfill$\square$

%---------------------------------------------------------------- 7
\exercise{7}{Let $A$ be a Noetherian ring. Then
\[
\dim A[x]=1+\dim A,
\]
and hence, by induction on $n$,
\[
\dim A[x_1,\dots,x_n]=n+\dim A.
\]
[Let $\fp$ be a prime ideal of height $m$ in $A$. Then there exist $a_1,\dots,a_m\in\fp$ such that $\fp$ is a minimal prime ideal belonging to the ideal $\fa=(a_1,\dots,a_m)$. By Exercise 7 of Chapter 4, $\fp[x]$ is a minimal prime ideal of $\fa[x]$ and therefore $\operatorname{height}\fp[x]\le m$. On the other hand, a chain of prime ideals $\fp_0\subset\fp_1\subset\cdots\subset\fp_m=\fp$ gives rise to a chain $\fp_0[x]\subset\cdots\subset\fp_m[x]=\fp[x]$, hence $\operatorname{height}\fp[x]\ge m$. Hence $\operatorname{height}\fp[x]=\operatorname{height}\fp$. Now use the argument of Exercise 6.]}
\pf
By Ex.~6, $\dim A[x]\ge1+\dim A$. For the reverse inequality it suffices to show:
\[
(\ast)\qquad\text{for every prime }Q\text{ of }A[x]\text{ with }\fp=Q\cap A:\quad\operatorname{ht}Q\le\operatorname{ht}\fp+1 ,
\]
because then $\dim A[x]=\sup_Q\operatorname{ht}Q\le\sup_\fp\operatorname{ht}\fp+1=\dim A+1$ (heights are finite by 11.12 since $A$ and, by the Hilbert basis theorem, $A[x]$ are Noetherian).

\emph{Reduction to $A$ local.} Let $T=A\setminus\fp$. Every prime $Q'\subseteq Q$ of $A[x]$ has $Q'\cap A\subseteq\fp$, so it does not meet $T$; by 3.11 and 3.13, $\operatorname{ht}Q=\operatorname{ht}(T^{-1}Q)$ in $T^{-1}(A[x])=A_\fp[x]$, and $T^{-1}Q\cap A_\fp=\fp A_\fp$. So we may replace $A$ by the Noetherian local ring $A_\fp$ and assume $A$ local with maximal ideal $\fp=\fm$, $d=\operatorname{ht}\fm=\dim A$ and $Q\cap A=\fm$.

\emph{The case $A$ local.} By 11.13 there are $a_1,\dots,a_d\in\fm$ with $\fa=(a_1,\dots,a_d)$ $\fm$-primary. By Ch.~4, Ex.~7(iii), $\fa[x]=\fa A[x]$ is $\fm[x]$-primary, so $r(\fa A[x])=\fm[x]$. Now $Q\supseteq\fm[x]$ and $Q/\fm[x]$ is a prime of $(A/\fm)[x]$ whose contraction to $A/\fm$ is $0$.
\begin{itemize}[leftmargin=*]
\item If $Q=\fm[x]$, then $Q$ is minimal over $\fa A[x]$ (its radical), an ideal generated by $d$ elements, so $\operatorname{ht}Q\le d$ by 11.16.
\item If $Q\supsetneq\fm[x]$, then $Q/\fm[x]$ is a non-zero prime of the PID $(A/\fm)[x]$, hence principal: $Q=\fm[x]+(f)$ for some $f\in A[x]$. The ideal $\fb=(a_1,\dots,a_d,f)$ is contained in $Q$, so $r(\fb)\subseteq Q$; and $r(\fb)\supseteq r(\fa A[x])=\fm[x]$ and $f\in r(\fb)$, so $r(\fb)\supseteq Q$. Thus $r(\fb)=Q$ and $Q$ is a minimal prime of $\fb$, which has $d+1$ generators; by 11.16, $\operatorname{ht}Q\le d+1$.
\end{itemize}
In either case $\operatorname{ht}Q\le d+1=\operatorname{ht}\fm+1$, proving $(\ast)$. Hence $\dim A[x]=1+\dim A$ (if $\dim A=\infty$, Ex.~6 gives $\dim A[x]=\infty$ as well).

\emph{Induction.} $A[x_1,\dots,x_n]=A[x_1,\dots,x_{n-1}][x_n]$ and $A[x_1,\dots,x_{n-1}]$ is Noetherian (7.5), so $\dim A[x_1,\dots,x_n]=1+\dim A[x_1,\dots,x_{n-1}]=n+\dim A$.

\emph{Remark (the hint's route).} The argument above also shows $\operatorname{ht}\fp[x]=\operatorname{ht}\fp$: the inequality $\ge$ comes from the chain $\fp_0[x]\subset\cdots\subset\fp[x]$, and $\le$ from the case $Q=\fm[x]$ treated above (after localizing at $\fp$, where $\fp[x]\mapsto\fm[x]$). \hfill$\square$

\section*{Notes}
There are no numbered in-text exercises in Chapter 11, and none of the 7 end-of-chapter exercises is only exposition. In Exercise 5 the phrase ``reformulate'' leaves the form of the answer open; the statement given (Theorem 11.1$'$) is the natural universal version, from which 11.1 for any additive function follows. Exercise 4 is a verification of the assertions in its statement (it cites Ch.~7, Ex.~9 for Noetherianity; that exercise is taken as known).

\end{document}
